\section{Involutions of the three-cusp group}

Throughout this part \(F=\widehat F(x,y)\), \(z=(xy)^{-1}\), and
\(\mathcal P\) denotes the subgroup of \(\operatorname{Out}(F)\)
preserving each named peripheral conjugacy class up to a common
unit power. All group homomorphisms are continuous. We use the
normalized lift
\[
 (\lambda,f):\quad x\longmapsto x^\lambda,\qquad
 y\longmapsto f^{-1}y^\lambda f,\qquad f\in F',
\]
where \(F'\) is the closed derived subgroup. This normalization is
a homomorphic section, as in the refereed LM4 D.32. Write
\(c=(-1,1)\). We use R1--R3 only with their stated scopes.

\statement{A.1}{Known}{Torsion in an acting group and its action on torsion are different assertions}
Lochak and Schneps (1997, Proposition 4(i)) prove
\(C_F(c)=1\); Proposition 4(ii) proves that \(c\) is
self-centralizing in \(\widehat{GT}\).
Schneps (1997, \S1.4(4)) and Lochak and Schneps (2006,
Question 1.4) ask whether all nonidentity torsion in
\(\widehat{GT}\) consists of conjugates of \(c\).
Harbater and Schneps (2000, Main Theorem(a), proved by
Propositions 3--4 and Theorem 2) identify
\(\widehat{GT}_0\) with the centralizer in \(\mathcal P\) of
the geometric \(S_3\).
Collas's genus-zero Theorems 3.3 and 3.10 identify prime-order
torsion in \(\widehat\Gamma_{0,[n]}\) with discrete torsion
and prove the \(\lambda\)-conjugacy action of \(\widehat{GT}\).
His genus-one Theorems 6.3 and 6.5 give the corresponding
statements with acting group \(\widehat{GR}\), isomorphic
to the subgroup denoted \(I\Gamma\) in that article;
replacing it there by all of \(\widehat{GT}\) is not
justified by those theorems.\cite{A_LS1997,A_Schneps1997,A_LS2006,A_HS2000,A_Collas2012a,A_Collas2012b}

\begin{proof}
For the centralizer comparison identify \(K(0,4)\) with
\(\widehat F(x,y)\). The natural \(S_4\)-action on its outer
automorphism group factors through the anharmonic \(S_3\).
Harbater--Schneps's quasi-special condition is precisely the
three peripheral power condition here; their symmetric
condition is commutation with that \(S_3\). Their proof
translates commutation with the transposition into the
two-cycle relation and commutation with the three-cycle into
the hexagon. The identification has the required source
and target, not merely an abstract isomorphism.

In Collas's formulas the variable of prime order lies in a
completed mapping class group; \((\lambda,f)\) is an arbitrary
element of the acting group. The conclusion
\(F(\gamma)=u\gamma^\lambda u^{-1}\) does not classify
finite-order elements \(F\) of the acting group. The subgroup
\(\widehat{GR}\) in the genus-one article is retained in its
Theorem 6.5. Neither action theorem therefore proves task A.8.
\end{proof}

The opened sources establish no literature antecedent for the
full order bound in R1. Searches included the authors' source
lists, the cited torsion papers, and recent work on profinite
curve-complex automorphisms and Galois--Teichm\"uller theory.
This records the search, not a priority claim or a proof that
no antecedent exists.

\statement{A.2}{Proved}{Relative pronilpotent levels recover the full peripheral group}
For an open characteristic subgroup \(U\) of \(F\), put
\[
 R_U=R_{\mathrm{nil}}(U),\quad F_U=F/R_U,\quad
 P_U=\operatorname{Out}^{\sharp}(F_U).
\]
For \(V\subset U\), the quotient \(F_V\to F_U\) induces
a continuous epimorphism \(P_V\to P_U\). There is an
isomorphism of profinite groups
\[
                 \mathcal P \xrightarrow{\sim}\varprojlim_U P_U.
\]
Two elements of \(\mathcal P\) are conjugate in \(\mathcal P\)
if and only if their images are conjugate in every \(P_U\).
This proves Lead 1.

\begin{proof}
Restricting finite nilpotent quotients of \(U\) to \(V\)
shows \(R_V\subset R_U\). Each \(R_U\) is characteristic in
\(F\). Since \(R_U\subset U\) and open characteristic
subgroups are cofinal in the finitely generated profinite
group \(F\), we have \(\bigcap_U R_U=1\).

We first prove surjectivity onto \(P_U\). Choose a representative
of a peripheral outer class with common multiplier \(\lambda\),
and a Galois element with cyclotomic multiplier \(\lambda\);
the cyclotomic character is onto
\(\widehat{\mathbb Z}^{\times}\). The latter acts on \(F_U\).
Composing with its inverse gives a multiplier-one peripheral
automorphism. R2 lifts this automorphism, with chosen
conjugators, to \(F\). Composing back proves surjectivity.
In particular every peripheral automorphism of \(F_V\)
lifts to \(F\), so preserves the image of \(R_U\).
This proves that the transition maps are defined on the
entire displayed outer groups. Their surjectivity follows
by lifting first to \(\mathcal P\).

Here the automorphism groups are compact: intersecting
kernels of all maps to finite groups of bounded order
gives a cofinal system of open fully invariant subgroups.
The conditions that a continuous endomorphism induces
bijections on these finite quotients are closed conditions
on its generator images. Consequently \(\operatorname{Aut}(F)\)
is profinite. Inner automorphisms form a compact subgroup.
The peripheral conditions are closed, since conjugators
and the common multiplier range over compact groups.
The same argument applies to the finitely generated
quotients \(F_U\).

If \([a]\in\mathcal P\) dies in every \(P_U\), the sets
\[
 \{g\in F:a(x)\equiv gxg^{-1},\
              a(y)\equiv gyg^{-1}\pmod {R_U}\}
\]
are nonempty compact sets with the finite intersection
property. Their intersection gives an actual inner
representative for \(a\), proving injectivity.
For a compatible family of classes in \(P_U\), their
nonempty compact inverse images in \(\mathcal P\) have
the same finite intersection property, proving surjectivity.
Finally apply this argument to the transporter sets
\(\{g:g_Ua_Ug_U^{-1}=b_U\}\) in \(\mathcal P\).
Surjectivity onto \(P_U\) makes these nonempty whenever
the two images are conjugate. Their intersection is the
required full conjugator.
\end{proof}

\statement{A.3}{Proved}{A central quotient is the exact obstruction to a peripheral lift}
Let \(N\) be a closed characteristic subgroup of \(F\), and
let \([N,F]\) denote the closed subgroup generated by
commutators between \(N\) and \(F\).
An automorphism \(a\) of \(F/N\), with specified conjugators
\(b_i\) satisfying
\[
 a(\bar x_i)=b_i\bar x_i b_i^{-1},\qquad
 (x_1,x_2,x_3)=(x,y,z),
\]
lifts to a peripheral automorphism of \(F\) with these
conjugators if and only if it lifts to a peripheral
\emph{automorphism} of \(F/[N,F]\) with conjugators above
the \(b_i\). This proves Lead 2.

\begin{proof}
Necessity follows from characteristicity. For sufficiency
let \(a^\#\) and \(b_i^\#\) be the data on \(F/[N,F]\).
We use the finite central-refinement count in the refereed
LM7\_PUSH D.50, equations (D.50.6)--(D.50.16).
For any actual quotient \(q_E:F\twoheadrightarrow E\)
and normal subgroup \(A\triangleleft E\), that count
constructs an actual quotient
\(q^\dagger:F\twoheadrightarrow D^\dagger\), with central
kernel over \(D=E/A\). A second epimorphism to
\(D^\dagger\) with specified peripheral conjugators has
the same number of epimorphism lifts to \(E\), with the
transported \(A\)-conjugacy classes on the boundary.
This counting assertion uses only normality of \(A\).
The hypothesis \([E,A]=A\) occurs in the separate
joint-image assertion (D.50.17), not used here.

Take \(A=q_E(N)\). The image \(q^\dagger(N)\) is central,
so \(q^\dagger\) kills \([N,F]\).
Thus \(q^\dagger a^\#\) is an epimorphism from the
same source to \(D^\dagger\), with conjugators the
images of \(b_i^\#\). The finite count gives an epimorphism
to \(E\) whose conjugators agree with the \(b_i\) in \(D\);
the reference count is positive since it includes \(q_E\).

In the compact space \(F^3\) of conjugators impose the
boundary product relation, generation in every finite
quotient, and agreement with each \(b_i\) in every finite
quotient of \(F/N\). A finite list of conditions lies in
one joint finite quotient \(E\), and the preceding
construction satisfies that list. Compactness gives
full conjugators satisfying all conditions. Their
endomorphism is surjective since its compact image
maps onto every finite quotient; Hopficity makes it
an automorphism. Its generator images induce exactly
\(a\) modulo \(N\).
\end{proof}

\statement{A.4}{Proved}{A degree-three coefficient imposed by the third cusp}
Over a field \(k\) of characteristic zero or greater
than three, use the free class-three Lie algebra on
\(X,Y\), and set
\[
 Z=[X,Y],\quad U_3=[X,Z],\quad V_3=[Y,Z],\quad
 x=\exp X,\quad y=\exp Y .
\]
If \(x\mapsto x^\mu,\ y\mapsto g^{-1}y^\mu g\), with
\(\mu\in k^\times\) and
\(\log g=tZ+\text{central terms}\), sends \(xy\)
to a conjugate of \((xy)^\mu\), then
\(t=(\mu^2-1)/24\).

\begin{proof}
The Baker--Campbell--Hausdorff formula gives
\[
 W=\log(xy)=X+Y+\tfrac12Z+\tfrac1{12}(U_3-V_3).
\]
The logarithm of the image is
\[
 \mu(X+Y)+\tfrac{\mu^2}{2}Z+
       \tfrac{\mu^3}{12}(U_3-V_3)+\mu tV_3 .
\]
Write the logarithm of a conjugator as
\(aX+bY+dZ+\text{central terms}\).
In \(\exp(\operatorname{ad}(aX+bY+dZ))(\mu W)\),
comparison of degree two gives \(m=a-b=(\mu-1)/2\).
The difference of the \(U_3\)- and \(V_3\)-coefficients,
divided by \(\mu\), then gives
\[
 \frac{\mu^2}{6}-t=\frac16+\frac m2+\frac{m^2}{2}
                  =\frac{3\mu^2+1}{24}.
\]
The result follows. The coefficient \(d\) and all
central coordinates cancel.
\end{proof}

\statement{A.5}{Refuted}{The specified conjugator in task A.3 need not preserve the third cusp}
There is an involution \(\alpha\in\widehat{GT}\), conjugate
to \(c\) by a Galois element, for which the unique solution
\(h\in F'\) of \(f_\alpha=h^{-1}\alpha(h)\) gives
\[
                 \psi(x)=x,\qquad\psi(y)=hyh^{-1}
\]
that does not preserve the conjugacy class of \(z\)
on \(F/R_{\mathrm{nil}}(F)\). The failure occurs in
the class-three exponent-seven quotient of order \(7^5\).

\begin{proof}
Choose \(\theta=(\mu,g)\) in the Galois image with
\(\mu\equiv2\pmod7\), possible by surjectivity of the
cyclotomic character modulo seven, and put
\(\alpha=\theta c\theta^{-1}\).
Normalization is a homomorphic section, so these
are actual automorphisms with that equality.
From \(\alpha\theta(y)=\theta(y)^{-1}\) we obtain
\[
 \alpha(y)=\alpha(g)g^{-1}y^{-1}g\alpha(g)^{-1},
 \qquad f_\alpha=g\alpha(g)^{-1}.
\]
Thus \(h=g^{-1}\) is a solution. Also
\(C_F(\alpha)=\theta(C_F(c))=1\) by
Lochak--Schneps Proposition 4(i). The quotient of two
solutions is fixed by \(\alpha\), so \(h\) is unique.

Let \(Q=\exp L\) for the free class-three Lie algebra
\(L\) on two generators over \(\mathbb F_7\).
It has order \(7^5\), exponent seven, and is generated
by \(x=\exp X,y=\exp Y\); hence it is a quotient of
\(F^{(7)}\) and of \(F/R_{\mathrm{nil}}(F)\).
The third-cusp condition on \(\theta\) and A.4 give
\[
       \operatorname{coeff}_Z(\log\bar g)
                    =\frac{2^2-1}{24}=1 .
\]
But \(\psi\) has multiplier one and the same conjugator
\(g\). Peripheral preservation for \(\psi\) would force
this coefficient to be zero by A.4, a contradiction.

The exact certificate \texttt{check\_A3\_class3.py}
realizes \(Q\) in the units of
\[
 \mathbb F_7\langle X,Y\rangle/(\text{words of length }4).
\]
It enumerates all \(16807\) elements; the conjugacy
orbit of \(xy\) has \(49\) elements. For
\(g=\exp[X,Y]\) it verifies
\[
 xg^{-1}yg\not\sim xy,\qquad
 x^2g^{-1}y^2g\sim(xy)^2 .
\]
The actual Galois conjugator differs only by central
degree-three coordinates relevant to this calculation,
which do not affect it. This certifies a failure in
an actual full profinite example; no order-preserving
lift of a finite involution is being assumed.
\end{proof}

The correct universal statement from R1 is that this
\(\psi\) intertwines \(\alpha,c\) on the maximal
pro-\(p\) quotients as automorphisms of free groups.
Peripheral preservation of this specific \(\psi\)
is an additional condition and is false.

\statement{A.6}{Proved}{The pure centralizer of standard conjugation has two elements}
\[
 C_{\mathcal P}(c)=N_{\mathcal P}(\langle c\rangle)=\{1,c\}.
\]
If \(\alpha\in\widehat{GT}_0\), every conjugator in
\(\mathcal P\) from \(c\) to \(\alpha\) belongs
to \(\widehat{GT}_0\).

\begin{proof}
A normalized \((\mu,g)\) commuting with \(c\) satisfies
\(c(g)=g\), by the normalized group law. A.1 gives
\(g=1\). In particular the multiplier-one centralizer
is trivial. Let \(\Sigma=S_3\) be the geometric
permutation group. If
\(\beta c\beta^{-1}=\alpha\) and \(\alpha,c\) commute
with \(\Sigma\), then
\(r_\sigma=\beta^{-1}\sigma\beta\sigma^{-1}\)
is pure, has multiplier one, and centralizes \(c\).
Thus \(r_\sigma=1\), and the Harbater--Schneps
identification in A.1 puts \(\beta\) in
\(\widehat{GT}_0\).

Apply this to a centralizing \((\mu,1)\).
The hexagon, with \(m=(\mu-1)/2\), reads
\((xy)^m=y^mx^m\). In the finite Heisenberg groups
with multiplication
\((a,b,d)(a',b',d')=(a+a',b+b',d+d'+ab')\),
the final coordinates of these two words are
\(m(m+1)/2\) and zero. Hence \(\mu^2=1\) in
\(\widehat{\mathbb Z}\). The actual automorphism
\((\mu,1)\) then has square one. R1 forces its
multiplier to be globally \(1\) or \(-1\), proving
the centralizer assertion. A normalizer of
\(\langle c\rangle\) centralizes its unique
nonidentity element.
\end{proof}

\statement{A.7}{Proved}{The multiplier-one cohomology set already has continuum many classes}
For \(G=\mathcal P,\widehat{GT}_0\), or \(\widehat{GT}\),
put \(H_G=\ker\lambda\), with \(C_2\) acting by \(c\).
There is a canonical closed injection of profinite spaces
\[
 \widehat{\mathbb Z}^{\times}/\{1,-1\}
 \longrightarrow H^1(C_2,H_G),\qquad
 \mu\longmapsto[\theta_\mu c\theta_\mu^{-1}c],
\]
where \(\theta_\mu\) is any element of multiplier \(\mu\).
The set \(H^1(C_2,H_G)\) has cardinality \(2^{\aleph_0}\).
For \(G=\widehat{GT}_0\), \(\Delta\) vanishes on this
closed family. All these \(H_G\)-conjugacy classes form
one \(G\)-conjugacy class.

\begin{proof}
A cocycle is \(a\in H_G\) with \(a\,c(a)=1\); it
corresponds to the involution \(ac\). Cohomology
identifies \(a\) and \(h^{-1}a\,c(h)\), exactly
conjugacy by \(H_G\). The Galois image makes
\(\lambda:G\to\widehat{\mathbb Z}^{\times}\) onto.
The element \(a_\mu=\theta_\mu c\theta_\mu^{-1}c\)
is a cocycle, and replacing \(\theta_\mu\) by a
different lift conjugates its involution by \(H_G\).

Two displayed involutions are \(H_G\)-conjugate
if and only if
\(\theta_\nu^{-1}h\theta_\mu\in C_G(c)=\{1,c\}\)
for some \(h\in H_G\), by A.6. Taking multipliers
gives \(\mu=\pm\nu\); these conditions are sufficient
as well. This proves the injection.
The cocycle space is compact, and its quotient by
the compact conjugating group is Hausdorff.
The orbit map from \(G\) factors continuously
through the multiplier quotient and then its
two signs, so the injection is continuous and closed.
Its domain has cardinality continuum. Conversely
\(G\) is contained in the metrizable profinite group
\(\operatorname{Out}(F)\), of cardinality at most
continuum. This proves equality.
Finally every represented involution is conjugate
to \(c\); the character \(\Delta\) is conjugacy
invariant and \(\Delta(c)=0\).
\end{proof}

\statement{A.8}{Proved}{Full conjugacy is an additional quotient of nonabelian cohomology}
With the notation of A.7, the full \(G\)-conjugacy
classes of nonidentity involutions are
\[
 \widehat{\mathbb Z}^{\times}\backslash H^1(C_2,H_G).
\]
The action of \(\mu\), using any lift \(\theta_\mu\), is
\[
 [a]\longmapsto[\,\theta_\mu a\,c(\theta_\mu^{-1})\,].
\]
For \(\mathcal N=\ker(\mathcal P\to P_F)\), the
\(\mathcal N\)-conjugacy classes of involutions with
image exactly \(c_F\) are \(H^1(C_2,\mathcal N)\).
Full \(\mathcal P\)-conjugacy on that fibre further
quotients by \(C_{P_F}(c_F)\), acting by lifted
conjugation, which need not preserve the distinguished
cohomology class.

\begin{proof}
By R1 every nonidentity involution has multiplier
\(-1\), so is uniquely \(ac\), \(a\in H_G\).
Its square-one condition is the cocycle equation.
Conjugation by \(\theta_\mu\) gives the formula.
Different lifts differ by \(H_G\), already divided
out in cohomology, proving the first assertion.

The same calculation for \(\mathcal P/\mathcal N\)
shows that involutions over the literal \(c_F\)
are \(ac\), \(a\in\mathcal N\). A conjugation
preserves this fibre exactly when its quotient
centralizes \(c_F\). Changing a lift changes
the action only by \(\mathcal N\)-conjugacy.
Such a lift need not centralize \(c\), explaining
the last qualification.
\end{proof}

\statement{A.9}{Proved}{Odd-prime conjugacy over a relative Frattini level}
Let \(U\) be open characteristic in \(F\), let \(p\)
be odd, put \(E_p=F/O^p(U)\), \(V=U/O^p(U)\),
and consider the peripheral outer subgroup preserving
\(V\). Its kernel on
\[
           E_p/\Phi(V)=F/\overline{U'U^p}
\]
is pro-\(p\). Two involutions with the same image
in this Frattini quotient are conjugate by the kernel.
This applies to the fibre over \(c\), where that
fibre is defined by descent to \(F/U\).

\begin{proof}
\(V\) is a nonabelian finitely generated free
pro-\(p\) group. Let \(J\) be the actual
automorphisms of \(E_p\) inducing identity on
\(E_p/\Phi(V)\). Restriction maps \(J\) into
the Frattini kernel of \(\operatorname{Aut}(V)\).
That kernel is pro-\(p\): on finite lower
\(p\)-central quotients its successive kernels
embed into homomorphism groups of elementary
abelian \(p\)-groups, and those quotients are cofinal.

Restriction is injective. If an element of \(J\)
fixes \(V\) pointwise, write \(a(g)=gu_g\) with
\(u_g\in\Phi(V)\). Comparing the conjugation
actions on \(V\) gives \(u_g\in C_V(V)=1\).
Thus \(J\) is pro-\(p\). If an outer automorphism
acts inner on \(E_p/\Phi(V)\), lift its inner
conjugator to \(E_p\) and correct it to an
element of \(J\). Consequently the outer
kernel is a continuous quotient of a closed
subgroup of \(J\), hence pro-\(p\).

The common image of the two involutions has
order two, since a pro-\(p\) group for odd \(p\)
has no involution. In its inverse image the
two cyclic subgroups are complements to the
normal pro-\(p\) kernel. Finite
Schur--Zassenhaus gives conjugators in every
finite quotient; compactness gives a conjugator
in the kernel itself.
\end{proof}

For arbitrary quotient automorphisms, preservation
of the image of a merely characteristic \(U\)
must be checked. It is included in the phrase
over the class of \(c\). For fully characteristic
\(U\) it is automatic: any quotient endomorphism
lifts to an endomorphism of the free profinite
group and preserves \(U\). Such \(U\) are cofinal.

\statement{A.10}{Proved}{Every negative pure pro-\(p\) involution is locally conjugate for odd \(p\)}
Let \(P^{(p)}\) be the pure peripheral outer group
of \(F^{(p)}\), with multiplier in
\(\mathbb Z_p^\times\), for odd \(p\).
Every involution of multiplier \(-1\) is
conjugate to \(c\) by a multiplier-one element.

\begin{proof}
The multiplier kernel \(H^{(p)}\) acts trivially
on abelianization, so lies in the pro-\(p\)
Frattini kernel used in A.9. The inverse image
of \(\{1,-1\}\) under the multiplier is
\(H^{(p)}\rtimes\langle c\rangle\).
Every negative involution defines a complement;
Schur--Zassenhaus gives the required conjugacy.
\end{proof}

This does not rescue the particular map of A.5.
Its local coboundary can be changed by a local
fixed element, whereas its full coboundary is
unique. The comparison of fixed subgroups
need not be surjective.

\statement{A.11}{Proved}{A peripheral reflection triangle is equivalent to full conjugacy}
For a nontrivial involution \(\alpha\in\mathcal P\)
let \(E=F\rtimes_\alpha C_2\). Then \(\alpha\)
is conjugate to \(c\) in \(\mathcal P\) exactly
when there are involutions \(t_0,t_1,t_2\)
representing its three section classes such that
\[
 a=t_1t_0,\qquad b=t_0t_2,\qquad d=t_2t_1
\]
generate conjugates of the respective inertia
groups of \(x,y,z\), and the \(t_i\) generate
\(E\) topologically. The last condition is
equivalent to generation in every \(E/R_U\).

\begin{proof}
For standard \(c\) take
\(t_0=c,t_1=xc,t_2=y^{-1}c\); their products
are \(x,y,z\) and they generate the extension.
A peripheral conjugator transports this triple.

Conversely the universal property gives a
surjection from the profinite coproduct
\(C_2^{(0)}\amalg C_2^{(1)}\amalg C_2^{(2)}\)
to \(E\), sending the factors to the named
involutions. Its parity kernel is free
profinite on \(a,b\). The restriction maps
onto \(F\): every even word in three
reflections is a word in the two edge products.
Both kernels are free profinite of rank two,
so Hopficity makes this restriction an
isomorphism. The equality \(abd=1\) is literal.
Abelianization makes the unit exponents in
the three peripheral conditions equal,
so the induced outer map belongs to \(\mathcal P\).
Conjugation by \(t_0\) inverts both \(a\) and \(b\).
Its outer action on \(F\) is \(\alpha\), since
\(t_0\) and the named splitting element have the
same image in \(E/F\). Hence the resulting outer
map conjugates \(c\) to \(\alpha\); no choice of
an actual splitting representative is needed.

Finally put \(H=\overline{\langle t_0,t_1,t_2\rangle}\).
If \(HR_U=E\) for all \(U\), then for every \(e\in E\)
the compact sets \(H\cap eR_U\) are nonempty
and have the finite intersection property.
Their intersection and \(\bigcap R_U=1\)
imply \(e\in H\). Thus \(H=E\).
\end{proof}

The transcription of cote 145, scan 23, lists
\(l_0=t_\infty t_1,\ l_1=t_0t_\infty,\
l_\infty=t_1t_0\). With ordinary multiplication
these satisfy \(l_0l_\infty l_1=1\).
The cyclic order or path-composition convention
must therefore be made explicit before
identifying them with \(x,y,z\).
We use \(abd=1\), and do not emend an unseen scan.

\statement{A.12}{Open}{Dyadic and mixed-prime involution obstructions}
A.9 does not determine the two-primary classes,
the gluing over \(U=F'F^3\) or \(U_{12}\), or
their \(\Delta\)-values. A.17 below supplies the
local framed equivariance of a parity-one construction,
but Lead 3 also needs an order-two full lift.
Neither R2 nor A.3 supplies that further equivariance.

For a surjection \(G\to Q\) with kernel \(L\),
a chosen lift \(t\) of an involution can be
corrected to an involution \(ut\), \(u\in L\),
exactly when
\[
                    u(tut^{-1})=t^{-2}.
\]
Surjectivity asserts the existence of \(t\),
not solvability of this nonabelian norm equation.
The map \(C_4\to C_2\) disproves that implication
in general, but is not a counterexample in
the peripheral groups. The residual task is
to solve the norm equations compatibly in
the actual relative-level kernels with the
specified frames and \(\Delta=1\), or to
prove they are insoluble.

\statement{A.13}{Open}{Parity on arbitrary involutions}
For an involution \(\tau=(-1,f)\) in
\(\widehat{GT}_0\), R3 gives
\(\Delta(\tau)=\delta((1,f))\).
It is zero for A.7's continuum family and
for every \(\widehat{GT}_0\)-conjugate of \(c\).
The cocycle relation alone gives only
\(2\Delta(a)=0\) in \(\mathbb F_2\).
We have neither an involution with
\(\Delta=1\) nor a proof of universal
vanishing, in \(\widehat{GT}_0\) or
\(\widehat{GT}\).

\statement{A.14}{Open}{The three full conjugacy assertions}
We have not decided whether every involution
of \(\mathcal P,\widehat{GT}_0,\widehat{GT}\)
is conjugate to \(c\) in its respective group.
A.2 detects conjugacy of a fixed pair on
all relative levels; A.11 identifies the
equivalent peripheral generation condition.
Neither supplies the compatible local
conjugators or reflection triangles whose
existence is required.

For \S31, folios 198--199 verso, scans
391--394, the full finite-subgroup conclusion
proved by R1 is that the only abstract
possibilities in \(\mathcal P\) are \(1,C_2\),
and \(C_2=\langle c\rangle\) occurs.
Exhaustion of its embedding conjugacy
classes is exactly the remaining assertion.

\statement{A.15}{Proved}{The pronilpotent lifting kernel also has continuum many cohomology classes}
Put \(Q=\mathcal P_F\), \(\mathcal N=\ker(\mathcal P\to Q)\),
and write \(c_Q\) for the image of \(c\). There is a canonical
closed injection
\[
 C_Q(c_Q)/\langle c_Q\rangle
 \longrightarrow H^1(C_2,\mathcal N),\qquad
 \bar\beta\longmapsto[\beta c\beta^{-1}c].
\]
Its image is exactly the \(\mathcal N\)-conjugacy classes in the
fibre over \(c_Q\) which are conjugate to \(c\) in \(\mathcal P\).
Both its domain and \(H^1(C_2,\mathcal N)\) have cardinality
\(2^{\aleph_0}\). Thus the large cohomology set does not itself
provide an additional full involution-conjugacy class.

\begin{proof}
The pronilpotent quotient is the product of its maximal
pro-\(p\) quotients. Every factor is its characteristic
Sylow subgroup; conversely, a tuple of continuous
automorphisms of the factors gives a continuous automorphism
of their product. Inner automorphisms and the three peripheral
conditions are coordinatewise. Therefore
\[
 Q=\prod_p\mathcal P^{(p)},\qquad
 C_Q(c_Q)=\prod_p C_{\mathcal P^{(p)}}(c).
\]
For odd \(p\), the multiplier map on the latter centralizer
is onto \(\mathbb Z_p^\times\). Indeed choose an element
\(\theta\) of the desired multiplier using the Galois image.
By A.10, some multiplier-one \(v\) conjugates
\(\theta c\theta^{-1}\) to \(c\); then \(v\theta\)
centralizes \(c\) and retains that multiplier.
Consequently the displayed centralizer, and its quotient
by the finite group \(\langle c_Q\rangle\), have at least
continuum many elements.

By A.2 every \(\bar\beta\in C_Q(c_Q)\) lifts to
\(\beta\in\mathcal P\). Its displayed commutator with
\(c\) lies in \(\mathcal N\) and is a cocycle.
Changing the lift changes its class only by
\(\mathcal N\)-conjugacy. Two lifts \(\beta_1,\beta_2\)
give the same class exactly when
\(\beta_2^{-1}n\beta_1\in C_{\mathcal P}(c)=\{1,c\}\)
for some \(n\in\mathcal N\). On taking the quotient,
this is exactly equality of the two cosets modulo
\(\langle c_Q\rangle\); the converse follows by lifting
that equality. This proves injectivity and the asserted
description of the image.
Compactness makes the injection continuous and closed,
as in A.7. All groups and orbit spaces here are compact
metrizable, so their cardinalities are at most continuum.
The preceding lower bound proves the asserted equalities.
\end{proof}

If the same construction is made in \(\widehat{GT}_0\),
using its actual image in any \(c\)-stable quotient,
every class in the connecting-map image has \(\Delta=0\).
This follows from conjugacy invariance of the character;
no surjectivity onto the entire peripheral quotient is
required for this observation.

\statement{A.16}{Proved}{A prescribed reflection splits every pro-two extension of the symmetric group}
Let \(1\to K\to E\xrightarrow{q}S_3\to1\) be an extension
of profinite groups with \(K\) pro-two. If a transposition
\(t\in S_3\) has an involution lift \(\widetilde t\in E\),
then there is a continuous section \(j:S_3\to E\) satisfying
\(j(t)=\widetilde t\).

\begin{proof}
For every finite \(\mathbb F_2[S_3]\)-module \(M\), restriction
\[
 H^n(S_3,M)\longrightarrow H^n(\langle t\rangle,M)
\]
is an isomorphism for \(n\geq1\). Indeed the double-coset
formula for restriction after corestriction has the identity
term and one term through the trivial intersection
\(\langle t\rangle\cap r\langle t\rangle r^{-1}=1\),
where \(r\) is a three-cycle. The second term is zero
in positive degrees. Corestriction after restriction is
multiplication by the odd index three, hence the identity
on these \(\mathbb F_2\)-vector spaces.

First suppose \(K\) is finite. Induct on its order.
When \(K\ne1\), the subgroup
\(A=\{z\in Z(K):z^2=1\}\) is a nonzero elementary
abelian characteristic subgroup of \(K\), hence normal
in \(E\). Induction gives a section of \(E/A\to S_3\)
with the prescribed image of \(\widetilde t\).
Pull back that section to an extension with kernel \(A\).
Its class in \(H^2(S_3,A)\) restricts to zero, since
\(\widetilde t\) is an involution. The proved injectivity
of restriction makes the extension split.
The difference between any splitting on \(\langle t\rangle\)
and the specified lift is an \(A\)-valued one-cocycle.
Surjectivity in degree one extends its cohomology class to
\(S_3\); changing the section by that cocycle and then by
an \(A\)-conjugation makes its value at \(t\) literally
\(\widetilde t\). This proves the finite assertion.

For general pro-two \(K\), take all open normal subgroups
\(L\) of \(E\) contained in \(K\). Each finite quotient
has such a section with the prescribed transposition.
The corresponding closed subsets of the compact space
of tuples \((e_g)_{g\in S_3}\in E^{S_3}\), imposing the
section identities modulo \(L\), are nonempty and have
the finite intersection property. Their intersection
gives a section with all identities literal.
\end{proof}

\statement{A.17}{Proved}{The parity-one relative frame can be chosen anti-invariant under conjugation}
In the notation of R2, put \(H_0=F'F^3\),
\(\Gamma_2=\widehat B/O^2(H_0)\), and let
\(d=\operatorname{Inn}(s)\circ c\), so that
\(d(s)=s\), \(d(r)=r^{-1}\), \(d(x)=y^{-1}\),
and \(d(y)=x^{-1}\). There is an actual element
\[
 b\in C_{H_0^{(2)}}(r),\qquad q(b)=[5I],\qquad d(b)=b^{-1}.
\]
Consequently the identity automorphism of
\(\widehat B/R_{\mathrm{nil}}(H_0)\), with \(s\)- and
\(\tau\)-frames one and \(r\)-frame equal to \(b\) at two
and one at all other primes, is a framed cocycle for
conjugation by \(c\). Its modulo-twelve frame parity is one.
This establishes local equivariance of these data; it does
not assert an equivariant full lift.

\begin{proof}
Use R2's actual joint quotient
\[
 J=\operatorname{im}\bigl(\widehat B\to Q\times D_0\bigr),
 \qquad Q=\mathrm{SL}_2(\mathbb Z/12)/\{\pm I\},\quad
 D_0=\widehat B/H_0.
\]
The map \(\Gamma_2\to J\) has pro-two kernel; both factors
and the joint quotient are preserved by \(d\).
The element \(z_0=([5I],1)\in J\) is central of order two
and is fixed by \(d\). Thus
\(S=\langle\bar r,z_0d\rangle\) in \(J\rtimes\langle d\rangle\)
is an actual copy of \(S_3\), with its named three-cycle
and transposition.

There is a particularly small literal lift of that
transposition as an involution. Set
\[
 C=[x,y],\qquad h=CxCx^{-1}\in F'\subseteq H_0.
\]
In R2's signed matrix coordinates modulo twelve,
\[
 q(h)=\begin{pmatrix}9&4\\8&9\end{pmatrix},\qquad
 D=\begin{pmatrix}0&1\\1&0\end{pmatrix},\qquad
 q(h)\bigl(Dq(h)D^{-1}\bigr)^{-1}=5I.
\]
Conjugation by \(D\) represents \(d\) on \(Q\).
The \(D_0\)-coordinate of \(h\) is one, so the actual
involution \(hdh^{-1}\) maps to \(z_0d\).
Apply A.16 to the preimage of \(S\) in
\(\Gamma_2\rtimes\langle d\rangle\), retaining this
involution. Its lifted three-cycle is conjugate to
the named \(r\) by the pro-two kernel, by the
coprime complement-conjugacy argument of A.9.
After that conjugation write the lifted transposition
as \(bd\). Its image in \(D_0\rtimes\langle d\rangle\)
is \(d\), so \(b\in H_0^{(2)}\), and its \(Q\)-coordinate
gives \(q(b)=[5I]\). Its order two gives \(b\,d(b)=1\).
Since it inverts the named \(r\), while \(d\) also
inverts \(r\), the element \(b\) centralizes \(r\).

All components of the prescribed relative frame have
image one in \(D_0\), so they give an actual element
of the relative fibre product. For a multiplier-one
normalized automorphism \(\Phi\) of \(\widehat B\),
write \(\Phi(s)=\beta s\beta^{-1}\) and
\(\Phi(r)=A rA^{-1}\). Direct substitution from
\(c(s)=s\), \(c(r)=sr^{-1}s\) gives the frames of
\(c\Phi c\):
\[
 \beta^c=c(\beta),\qquad
 A^c=c(\beta)s\,c(\beta)^{-1}c(A)s.
\]
For the specified local data \(\beta=1\), the second
formula is \(A^c=d(A)=A^{-1}\). These are exactly
the frames of the inverse of the identity framed
automorphism. The parity is the same as in R2's E.58:
the signed \(r\)-frame is \(5I\), whose modulo-three
value is \(-I\), and the \(s\)-frame is one.
The independently executed certificate
\texttt{check\_A17\_frame.py} verifies the literal
matrix word and the covariance identities just used.
\end{proof}

The remaining full obstruction can now be placed in a
smaller, specified kernel. Let \(L\) consist of the
multiplier-one elements whose normalized action is
identity on \(\widehat B/R_{\mathrm{nil}}(H_0)\) and whose
two pure frames both have image one there. It is the
kernel of the framed action, is stable under \(c\), and
\(\Delta(L)=0\). R2 supplies a full \(\sigma\) lifting
A.17's data, with \(\Delta(\sigma)=1\). Put \(a=\sigma c\).
Then \(a^2\in L\). A parity-one involution in this
specified fibre exists exactly when some \(v\in L\)
solves
\[
                   v\,(ava^{-1})=a^{-2}.
\]
The local construction proves that the right-hand side
lies in this framed kernel. Solvability of this last full
norm equation remains open.

\statement{A.18}{Refuted}{R2 does not admit an unrestricted order-preserving refinement}
There is an actual involution of
\(\operatorname{Out}^{\sharp}(F/R_{\mathrm{nil}}(F))\)
which lifts to a full peripheral automorphism by R2 but
has no involution lift in \(\mathcal P\).

\begin{proof}
Use the actual product
\(F/R_{\mathrm{nil}}(F)=\prod_pF^{(p)}\).
Define \(\bar a\) to be standard conjugation \(c\)
on \(F^{(3)}\) and identity on every other factor.
This is a continuous automorphism of order two,
preserving each of the three peripheral conjugacy
classes up to the common unit
\[
 \lambda_3=-1,\qquad \lambda_p=1\quad(p\ne3).
\]
Its outer class is nonidentity, since its action on
the three-primary abelianization is \(-I\).
By A.2, which applies R2 after Galois translation
of the multiplier, this class has a lift in
\(\mathcal P\). Every lift has the displayed
multiplier, because the product quotient retains
the entire abelianization \(\widehat{\mathbb Z}^2\).
By R1, a nonidentity full involution must instead
have multiplier \(-1\) at every prime. Hence no
lift of \(\bar a\) can be an involution.
\end{proof}

This counterexample is in the actual peripheral level
groups, rather than in an unrelated finite extension.
It excludes a blanket equivariant strengthening of
the central-refinement count. A.17's remaining fibre
has multiplier globally \(-1\), so this orientation
obstruction does not decide that fibre.

\statement{A.19}{Proved}{The square obstruction at an abelian refinement and the character-lifting test}
Suppose \(a\) is a lift of a quotient involution,
\(a^2\in L\), and \(L_0\triangleleft\langle L,a\rangle\)
is contained in \(L\), with \(M=L/L_0\) finite abelian.
Conjugation by \(a\) makes \(M\) a \(C_2\)-module.
An involution exists above that quotient element in
\(\langle L,a\rangle/L_0\) exactly when
\[
 [a^2]=0\quad\hbox{in}\quad
 H^2(C_2,M)=M^{a}/(1+a)M.
\]
This obstruction vanishes when \(M\) has odd order.
For elementary abelian two-primary \(M\), its target is
\(\ker(a+1)/\operatorname{im}(a+1)\).

If \(\Delta\) has a continuous lift to a character
\(\widehat{GT}_0\to\mathbb Z/4\), then it vanishes on
every involution. The same conclusion follows from
a continuous lift of \(\delta\) on \(\mathcal H\)
which is invariant under conjugation by \(c\).
Existence of either lift is not established here.

\begin{proof}
The square of the conjugation action is conjugation
by \(a^2\in L\), which is trivial on \(M\).
Also \(a^2\) is invariant under conjugation by \(a\).
Replacing \(a\) by \(va\), with \(v\in L\), changes
its square in additive notation by \((1+a)\bar v\).
This proves the quotient criterion, which is the
usual cyclic two-cocycle description of \(H^2\).
For an invariant \(u\) of odd order, \(u/2\) is
defined and \((1+a)(u/2)=u\), proving vanishing.
In characteristic two the invariants are
\(\ker(a+1)\), giving the last formula.

A character to \(\mathbb Z/4\) sends an involution
to zero or two; both reduce to zero modulo two.
Alternatively, if \(b\in\mathcal H\) is a cocycle
for \(c\) and \(\theta\) is a \(c\)-invariant
lift of \(\delta\), then
\(0=\theta(bc(b))=2\theta(b)\), so
\(\delta(b)=0\).
\end{proof}

For the actual \(a=\sigma c\) of A.17, any proposed
global character lift \(\Theta\) would necessarily satisfy
\(\Theta(a^2)=2\), although \(a^2\in L\) and
\(\Delta(L)=0\). Thus such a character must detect the
residual square inside the framed kernel; it cannot be
inferred from the local frames alone. Conversely A.17's
local parity-one framed involution rules out a character
lift to \(\mathbb Z/4\) on that local framed group itself.

\statement{A.20}{Proved}{The locally glued symmetric subgroup does not lift in the standard extension}
Let \(H_0\), \(d\), and \(b\) be as in A.17. In the standard
relative extension
\[
 \bigl(\widehat B/R_{\mathrm{nil}}(H_0)\bigr)
                  \rtimes\langle d\rangle
\]
there is a subgroup \(S_{\mathrm{loc}}\simeq S_3\), generated
by the named \(r\) and a transposition with component \(bd\)
at two and \(d\) at every odd prime. No finite subgroup
of the standard full extension
\(\widehat B\rtimes\langle d\rangle\) maps onto
\(S_{\mathrm{loc}}\). The obstruction is visible already in
\(Q\rtimes\langle d\rangle\), with \(|Q|=576\).

\begin{proof}
All local components of the transposition have the same
image \(d\) in \(D_0\rtimes\langle d\rangle\).
They therefore form an actual element of the relative fibre
product. A.17 gives its order two and its inversion of \(r\).
Their images in \(D_0\rtimes\langle d\rangle\) generate
\(S_3\), so the resulting subgroup has order exactly six.

In the full free product \(\widehat B=C_2\amalg C_3\),
every element of order three is conjugate into the named
\(C_3\), and that factor is malnormal. These are
Herfort--Ribes Theorems A' and B'.\cite{A_HR1985}
In particular \(N_{\widehat B}(\langle r\rangle)=\langle r\rangle\).
Since \(d\) inverts \(r\),
\[
 N_{\widehat B\rtimes\langle d\rangle}(\langle r\rangle)
                         =\langle r,d\rangle\simeq S_3.
\]
Every full subgroup isomorphic to \(S_3\) intersects
\(\widehat B\) in its unique subgroup of order three:
it cannot be contained in \(\widehat B\), whose finite
subgroups conjugate into its factors by the same paper's
Theorem 1. Conjugating that order-three subgroup into
\(\langle r\rangle\) proves that every full \(S_3\)
is conjugate to \(\langle r,d\rangle\).

The image of \(S_{\mathrm{loc}}\) in the actual finite
quotient \(Q\rtimes\langle d\rangle\) is
\[
                   \langle R,[5I]d\rangle .
\]
It is not conjugate there to \(\langle R,d\rangle\).
Indeed a conjugator may be taken in \(Q\), since \(d\)
normalizes the latter subgroup, and it must normalize
its unique order-three subgroup. Directly solving the
two-by-two equations in \(\mathrm{SL}_2(\mathbb F_3)\)
shows that no determinant-one matrix conjugates \(R\)
to \(R^{-1}\), and its centralizer is \(\{\pm R^j\}\).
The centralizer modulo four is also \(\{\pm R^j\}\).
Traces exclude a negative sign in the projective
commutation equation. Thus a representative of a
normalizing element has coordinates
\(\varepsilon_4R^i\) modulo four and
\(\varepsilon_3R^j\) modulo three.
Conjugation by \(D\) fixes both signs and inverts \(R\).
Its \(d\)-norm consequently has positive signs in both
coordinates. It cannot lie in \([5I]\langle R\rangle\):
one global choice of representative has negative sign
modulo three, and the other has negative sign modulo four.
This proves the nonconjugacy and hence the asserted
nonlifting of an \(S_3\) subgroup.

For completeness, a finite subgroup mapping onto
\(S_{\mathrm{loc}}\) has trivial kernel: the kernel lies
in the torsion-free free profinite group \(H_0\).
It would therefore be an \(S_3\), already excluded.
The exact exhaustive certificate
\texttt{check\_A20\_standard\_extension.py} generates all
\(576\) elements of \(Q\), obtains normalizer and
centralizer orders \(18\), lists all nine \(d\)-norm
values and all three elements of the target coset, and
checks that the lists are disjoint. It independently
enumerates the two signed centralizers used above.
\end{proof}

A.20 blocks lifting this subgroup inside the \emph{standard}
extension. A hypothetical new involution of
\(\widehat{GT}_0\) would change the action defining the
full extension. No obstruction in that unknown twisted
extension is inferred from A.20.

\statement{A.21}{Refuted}{Real projectivity and a rank-two open free subgroup do not imply finite graph recognition}
There exists a finitely generated real-projective profinite
group \(E_*\) with an open normal subgroup isomorphic to
\(\widehat F_2\), quotient \(C_2\), and self-centralizing
involutions, which is not the profinite fundamental group
of a finite graph of finite groups. It has one conjugacy
class of involutions, so it does not satisfy R1's
three-section conclusion.

\begin{proof}
Put \(A=\prod_{p\text{ odd}}\mathbb Z_p\) and define
\[
 \Gamma_*=A\rtimes\mathbb Z_2,
 \qquad u\cdot a=(-1)^{u\bmod2}a.
\]
All groups and actions here are profinite and continuous.
The kernel of the parity map \(\Gamma_*\to C_2\) is
\(A\times2\mathbb Z_2\simeq\widehat{\mathbb Z}\).
The group \(\Gamma_*\) is torsion-free: a finite-order
element projects trivially to \(\mathbb Z_2\), and then
belongs to the torsion-free group \(A\).
Set
\[
 E_*=C_2\amalg\Gamma_*,\qquad
 \varepsilon:E_*\longrightarrow C_2,
\]
where \(\varepsilon\) is identity on the first factor
and parity on the second. In the open-subgroup Kurosh
decomposition of \(\ker\varepsilon\), the two vertex
groups are \(1\) and \(\widehat{\mathbb Z}\), and the
quotient graph has two vertices joined by two edges.
Consequently
\[
 \ker\varepsilon\simeq
 \widehat{\mathbb Z}\amalg\widehat{\mathbb Z}
 \simeq\widehat F_2.
\]
The open Kurosh theorem in precisely this full-profinite
form is recalled in Herfort--Ribes, Section~2.
Their Theorems A', B', and 1 show that all nonidentity
torsion lies in conjugates of the single \(C_2\)
factor, and that each such involution has centralizer
exactly that factor.\cite{A_HR1985}
Haran's characterization of real projectivity therefore
applies: an index-at-most-two subgroup of cohomological
dimension at most one, together with self-centralizing
involutions, suffices. We use its exact formulation in
P\'al--Quick, Theorem~2.1.\cite{A_PalQuickFormal2025}

Suppose \(E_*\) were the profinite fundamental group of
a finite graph of finite groups. Its vertex groups have
order at most two. Collapse each nonloop edge whose
nontrivial edge group is equal to an endpoint group.
In the resulting reduced graph, a loop with edge group
\(C_2\) would make its stable letter centralize that
\(C_2\); the stable letter has infinite order, as is
seen by killing the vertex groups and mapping that
loop to \(\widehat{\mathbb Z}\). This contradicts the
centralizer just proved. All remaining edges therefore
have trivial group. We obtain
\[
 E_*\simeq\bigl(\mathop{\amalg}_{j=1}^{k}C_2\bigr)
                   \amalg\widehat F_r .
\]
The single torsion class forces \(k=1\). The
torsion-free index-two subgroup of this free product
has rank \(2r+k-1\), by its two-sheeted graph or by
open Kurosh. The rank two already computed forces
\(r=1\).

On the other hand \([\Gamma_*,\Gamma_*]=A\), since
multiplication by two is invertible on \(A\).
Thus
\[
 E_*^{\mathrm{ab}}=C_2\times\mathbb Z_2,
\]
whereas \((C_2\amalg\widehat{\mathbb Z})^{\mathrm{ab}}
=C_2\times\widehat{\mathbb Z}\). The latter has
nontrivial odd-primary quotients and the former does not.
This contradiction proves the claim.
\end{proof}

The construction uses the torsion-free virtually
procyclic mechanism already present in Haran's
counterexamples to full-profinite virtual freeness
recognition.\cite{A_Haran1985}
It identifies a genuine limitation of that method;
it is not a new peripheral involution. For the groups
in A.14 the three-section hypothesis, and ultimately
the marked cusp products in A.11, still have to be used.

\statement{A.22}{Proved}{The three distinguished reflections normally generate, and a proper generated subgroup has infinite index}
Let \(\alpha\) and \(h\) be as in R1, and put
\(E=F\rtimes_\alpha\langle t\rangle\), \(b=hyh^{-1}\).
The three involutions
\[
       t_0=t,\qquad t_1=xt,\qquad t_2=b^{-1}t
\]
represent the three section classes and normally
generate \(E\). If their closed generated subgroup
\(J\) is proper, then \([E:J]\) is infinite.
Normal generation alone does not supply the full
generation required in A.11.

\begin{proof}
The equality \(f=h^{-1}\alpha(h)\) gives
\(\alpha(b)=b^{-1}\), so all three displayed
elements have order two. In
\(E^{\mathrm{ab}}\simeq C_2^3\), their images are
\(t,x+t,y+t\); these are distinct and hence the
elements are pairwise nonconjugate. R1 states that
there are exactly three classes. Quotienting \(E\)
by their closed normal closure kills \(t\), then
\(x=t_1t\), then \(b^{-1}=t_2t\). Since \(b\) is a
conjugate of \(y\), it kills \(y\), and the quotient
is trivial.

There is a continuous epimorphism
\(C_2\amalg C_2\amalg C_2\to J\) sending the three
factors to the displayed involutions. Its parity
kernel is free profinite of rank two, so \(J\cap F\)
is generated by at most two elements. If \(J\)
had finite index in \(E\), then \(J\cap F\) would
be open in \(F\) and the profinite Schreier formula
would give
\[
 d(J\cap F)=1+[F:J\cap F].
\]
It follows that \(J\cap F=F\), and \(t\in J\)
then gives \(J=E\). This proves the assertion.
\end{proof}
